% Generated draft: requires mathematical and source review.
\documentclass[11pt]{article}
\usepackage[margin=25mm]{geometry}
\usepackage{fontspec}
\setmainfont{Latin Modern Roman}
\usepackage{amsmath,amssymb,mathtools,unicode-math}
\setmathfont{Latin Modern Math}
\usepackage{enumitem,needspace,xurl,hyperref}
\setlength{\parindent}{0pt}
\setlength{\parskip}{5pt}
\setlength{\emergencystretch}{3em}
\newcommand{\N}{\mathbb{N}}
\newcommand{\Z}{\mathbb{Z}}
\newcommand{\Q}{\mathbb{Q}}
\newcommand{\R}{\mathbb{R}}
\newcommand{\C}{\mathbb{C}}
\newcommand{\F}{\mathbb{F}}
\newcommand{\A}{\mathbb{A}}
\newcommand{\PP}{\mathbb{P}}
\newcommand{\E}{\mathbb{E}}
\newcommand{\eps}{\varepsilon}
\newcommand{\dd}{\,\mathrm{d}}
\begin{document}
\section*{UnsolvedMath ID: 800165}
\textbf{Source problem code:} catalogue entry
\subsection{Definitions and mathematical statement}
Let \(\mathcal H\) denote the intended family of holomorphic Hecke eigenforms, and let \(k_f\) be the weight of \(f\in\mathcal H\). Write \(L(s,f)\) for the associated \(L\)-function, with the normalization implicit in the source's designation of \(s=\tfrac12\) as the central point, and \(L(s,\operatorname{ad}f)\) for its adjoint \(L\)-function. For real \(T>1\) and \(r\in\{3,4\}\), put
\[
M_r(T)=\sum_{\substack{f\in\mathcal H\\ k_f\leq T}}
\frac{L(\tfrac12,f)^r}{L(1,\operatorname{ad}f)}.
\]
The family must be specified so that this sum is well defined, including finite multiplicities and defined, nonzero denominators. The supplied record does not specify the level, nebentypus, whether only cusp forms are included, the admissible weights, or the convention for counting eigenforms. These unresolved qualifications are essential to a precise statement; no particular choice is imposed here.

\textbf{Requested asymptotics, subject to specifying the intended family.} Prove that there exist fixed polynomials \(P_3\) and \(P_6\), independent of \(T\) and \(\varepsilon\), with
\[
\deg P_3=3,\qquad \deg P_6=6,
\]
such that, for every \(\varepsilon>0\), as \(T\to\infty\),
\[
M_3(T)=T^2P_3(\log T)+O_{\varepsilon}(T^{1+\varepsilon})
\]
and
\[
M_4(T)=T^2P_6(\log T)+O_{\varepsilon}(T^{4/3+\varepsilon}).
\]
The implied constants may depend on \(\varepsilon\) and the fixed family conventions, but not on \(T\). The coefficients of the polynomials are not supplied.

The source additionally asks for the best possible error terms and expresses the expectation that the displayed bounds are accessible using current technology. It does not specify a separate optimality criterion or supply lower bounds establishing that these exponents are optimal. It reports analogous results for Hecke--Maass cusp forms: the fourth moment is attributed to Ivic and the third moment to Zhi Qi, whose result is described as having appeared ``earlier this year'' in the 2024 source context. These are supplied contextual claims, not independently verified assertions here.
\subsection{Short English statement}
For the intended family of holomorphic Hecke eigenforms of weight at most \(T\), establish harmonic third and fourth central \(L\)-value moment asymptotics with polynomial main terms of degrees \(3\) and \(6\), and errors \(O_{\varepsilon}(T^{1+\varepsilon})\) and \(O_{\varepsilon}(T^{4/3+\varepsilon})\), respectively. The family and counting conventions require clarification. The source also asks for best possible errors without defining optimality.
\subsection{Sources}
\begin{itemize}
\item Local UnsolvedMath registry and supplied complete source record: permanent ID \(800165\), exact source problem code \texttt{catalogue entry}, upstream revision \texttt{d371ee917ac9647fbb06cd6fc529deec6f3dda82}. The supplied registry status is \texttt{open}; this is a source label, not an independently verified assessment.
\item Hamieh, Kadiri, Ng, \emph{Moments of L-functions Problem List} (2024), Problem 15, \url{https://arxiv.org/abs/2405.17800}. The source attributes the problem to Peter Humphries, Virginia; the registry's proposer field lists Hamieh, Kadiri, Ng. The supplied provenance records access on 2026-10-16 and extraction by \texttt{arxiv\_latex}; no source visit is claimed here.
\item The supplied statement attributes an analogous Hecke--Maass fourth-moment result to Ivic and a third-moment result to Zhi Qi, described as proved ``earlier this year'' in the 2024 source context. No separate bibliographic references or links for these results are supplied.
\end{itemize}
\end{document}
